\chapter{多智能体协同}

\section{为什么不做“自由辩论”}

榜题在创新性条款中点名要求“交叉验证与辩论”机制。
本方案\emph{没有}让两个模型来回争辩三轮——那样既消耗大量令牌，
又无法测量，对准确率也没有帮助。这里实现的是\emph{断言粒度的对抗与仲裁}。

第一步是两位专家独立起草。关键在于“独立”：专家甲按知识点硬过滤取前若干条切片，
视角窄而准；专家乙全库检索后按知识点重排，视角宽而全。
\emph{视角不同，犯的错才不同}。两个智能体如果看到同一份上下文，
等于同一个模型问了两遍，没有任何信息增益。

第二步是对齐。按相似度将两份断言配对，分为双方都提出、仅一方提出、双方冲突三类。

第三步是仲裁。裁判不投票，也不看谁表述得更漂亮，只回到知识库比对证据。
证据分差不足时两条都弃用。

\section{判定“一致”的三个条件}

\begin{keynote}{整个机制里最关键的一行}
判定一致必须同时满足三条：文本高度相似、出处相同、\emph{数值集合相同}。

第三条是被单元测试逼出来的。原先的实现只看前两条，
结果专家甲说“限速 250\unit{mm/s}”、专家乙说“限速 600\unit{mm/s}”，
文本相似度高达 $0.95$、引用出处又完全相同，被判成“双方印证”直接放行。
\emph{数值篡改恰恰是文本最相似的一类冲突}，只看相似度等于把最危险的情形
当成了最可信的情形，方向完全反了。
\end{keynote}

对应地，仲裁分两层：数值集合不同时，以“谁的数字在知识库里”定胜负，黑白分明；
其余情况才比较整体证据分。

\section{辩论机制的实测效果}

关闭审核环节、只观察辩论环节的分层统计，
在同一批 9168 条断言上的结果见表 \ref{tab:consensus}。

\begin{table}[htbp]\centering\small
\caption{按共识等级分层的无依据率}\label{tab:consensus}
\begin{tabular}{@{}p{3.4cm}p{2.4cm}p{2.6cm}p{2.4cm}@{}}
\toprule
共识等级 & 断言数 & 无依据条数 & 无依据率 \\
\midrule
双方独立印证 & 5001 & 27 & 0.54\% \\
仅单方提出 & 4167 & 786 & 18.86\% \\
\bottomrule
\end{tabular}
\end{table}

两者相差约三十五倍。这说明“两位专家是否独立提出同一条”本身
就是一个强幻觉信号，交叉验证不是走过场：
一条被编造出来的断言，另一位采用不同检索视角的专家几乎不会编出同一句。

\begin{keynote}{必须写进报告的限制}
上面这组数是\emph{离线桩}跑出来的。桩按种子给两位专家注入不同的假事实，
“独立编造不会撞车”在桩里部分是构造出来的。
真实模型下这个倍数会变，\emph{接上真模型后必须重测}，不要把三十五倍当作结论。
\end{keynote}

\bibliographystyle{gbt7714-2005}
\bibliography{refs}
